\chapter{SISTEMAS DE DETECÇÃO DE INTRUSÃO}

Este capítulo apresenta os fundamentos dos Sistemas de Detecção de Intrusão, estabelecendo o referencial conceitual que embasa as escolhas metodológicas deste trabalho. São abordadas as classificações dos IDS por origem dos dados monitorados, as duas abordagens de detecção, por assinatura e por anomalia, e os principais desafios operacionais que motivam a adoção de técnicas de aprendizado de máquina como alternativa às defesas baseadas em regras estáticas.

\section{Conceitos fundamentais de IDS}

Um Sistema de Detecção de Intrusão (\textit{Intrusion Detection System} - IDS) é uma ferramenta de segurança projetada para monitorar, analisar e alertar sobre atividades suspeitas ou maliciosas em redes de computadores ou sistemas de informação. Essas atividades incluem tentativas de acesso não autorizado, exploração de vulnerabilidades, ataques de \textit{malware} e outras ameaças à segurança \cite{santos2023ercemapi}. O controle exercido pelo IDS ocorre sobre eventos do sistema operacional ou sobre o tráfego de rede, buscando possíveis violações de normas e políticas de segurança da informação \cite{liborio2015}.

O conceito de detecção de intrusão foi formalizado por Dorothy Denning em 1987, no trabalho \textit{An Intrusion-Detection Model}. Nesse trabalho, Denning estabeleceu que atividades intrusivas geram rastros anômalos que desviam dos padrões habituais de uso dos sistemas, tornando viável a análise contínua de registros de auditoria para distinguir atividades maliciosas de operações legítimas \cite{ids_introducao}. Na década de 1990, Todd Heberlein, da Universidade da Califórnia em Davis, introduziu a detecção de intrusão baseada em rede com o desenvolvimento do \textit{Network Security Monitor} (NSM), o primeiro sistema voltado especificamente para a análise do tráfego de rede. O NSM foi empregado em instalações governamentais e gerou um crescimento significativo de interesse e investimento na área \cite{ids_introducao}.

Os IDSs são classificados de acordo com a origem dos dados monitorados. Os IDSs baseados em rede (\textit{Network-based IDS}, NIDS) são implantados em pontos estratégicos da infraestrutura e examinam os pacotes à medida que trafegam pela rede, sem exigir a instalação de \textit{software} em cada \textit{host}. Essa característica torna os NIDS mais escaláveis, capazes de monitorar dezenas ou centenas de dispositivos por meio de um único sensor \cite{santos2023ercemapi}. Os IDSs baseados em \textit{host} (\textit{Host-based IDS}, HIDS) operam em nível de sistema operacional, analisando uso de memória, conexões ativas, chamadas de sistema e arquivos de \textit{log} de uma única máquina. Embora ofereçam visibilidade mais profunda sobre ataques que comprometem o sistema alvo, demandam a instalação e manutenção de agentes em cada máquina monitorada \cite{liborio2015}. Sistemas híbridos combinam as duas abordagens, aplicando primeiro a detecção por assinatura e, na ausência de correspondência, recorrendo à detecção por anomalia \cite{liborio2015}.

\section{Abordagens de detecção}

A detecção por assinatura, também chamada de detecção por uso indevido, opera comparando o tráfego ou os eventos do sistema com um banco de dados de padrões conhecidos de ataques. O funcionamento é análogo ao de um software antivírus: quando os dados analisados correspondem a uma assinatura cadastrada, o sistema gera um alerta \cite{liborio2015}. A principal vantagem dessa abordagem é a baixa taxa de falsos positivos, pois o alerta é produzido apenas quando há correspondência direta com um ataque documentado. A limitação fundamental é que o sistema permanece sem resposta diante de ataques desconhecidos ou variações de ataques conhecidos que não estejam representados na base. A base de assinaturas também precisa ser atualizada de forma contínua para manter a efetividade \cite{ravindran2025}.

A detecção baseada em anomalias parte da premissa oposta: em vez de reconhecer ataques conhecidos, o sistema constrói um perfil do comportamento normal da rede ou do host e sinaliza qualquer desvio significativo desse perfil. Para isso, o IDS passa por uma fase de treinamento na qual estabelece uma linha de base (\textit{baseline}) com métricas do tráfego legítimo. Padrões como grande fluxo de pacotes UDP, crescimento anormal de tráfego ICMP ou excesso de requisições \textit{web} são exemplos de anomalias detectáveis por esse método \cite{liborio2015}. A principal vantagem é a capacidade de identificar ataques inéditos, cujas assinaturas ainda não foram catalogadas. A contrapartida é uma taxa mais elevada de falsos positivos, pois variações legítimas no tráfego também podem ser interpretadas como anomalias \cite{santos2023ercemapi}.

A adoção de técnicas de aprendizado de máquina permitiu enriquecer ambas as abordagens de detecção, mas com papéis distintos. No paradigma supervisionado, algoritmos como \textit{Random Forest}, \textit{Support Vector Machine} (SVM), $k$-\textit{Nearest Neighbors} ($k$NN), redes neurais artificiais (ANN) e \textit{XGBoost} são treinados com dados rotulados contendo exemplos explícitos de tráfego benigno e de cada classe de ataque; o modelo aprende a classificar novos fluxos entre essas categorias predefinidas \cite{santos2023ercemapi}. Trata-se, portanto, de uma abordagem de classificação supervisionada, e não de detecção de anomalias em sentido estrito: o modelo só reconhece classes que estiveram presentes no treinamento. No paradigma não supervisionado, modelos como \textit{Autoencoders} são treinados exclusivamente com tráfego legítimo e sinalizam como ataque qualquer fluxo cujo comportamento se desvie significativamente do padrão aprendido, sem necessidade de rótulos de ataque. Essa é a forma de detecção mais próxima do conceito original de detecção por anomalia.

\section{Desafios atuais de detecção em redes}

Um dos desafios mais documentados na área é a taxa de falsos positivos. Em ambientes reais, variações legítimas no tráfego de rede, como campanhas de atualização de sistemas ou picos de acesso em horários específicos, podem disparar alertas incorretos em sistemas baseados em anomalias. O volume de alertas gerado sobrecarrega as equipes de operações de segurança e reduz a confiança no sistema \cite{santos2023ercemapi}.

A detecção de ataques inéditos representa outra limitação estrutural, em especial para sistemas baseados em assinatura. Novas variantes de ataques são desenvolvidas de forma contínua, e a janela entre a descoberta de uma vulnerabilidade e a criação da assinatura correspondente deixa os sistemas expostos. Técnicas de evasão, como ofuscação de código e modificação da estrutura do ataque, são utilizadas para contornar assinaturas existentes \cite{ravindran2025}.

A qualidade e a representatividade dos conjuntos de dados de treinamento são determinantes para o desempenho dos modelos de aprendizado de máquina. \textit{Datasets} com desequilíbrio severo de classes, nos quais o tráfego benigno supera o malicioso em várias ordens de magnitude, dificultam o aprendizado de padrões de ataque menos frequentes. Modelos treinados em um ambiente específico podem não generalizar para redes com características de tráfego distintas, um problema amplamente investigado na literatura \cite{neto_gomes}.
